% Coverage: Chapter 8, Section 8.4, physical PDF pp. 373--376
% (printed pp. 350--353; p. 376 only through the line before the Section 8.5 heading).
% Semantic range: Section 8.4 heading through "because a^mu has been defined to
% have a^0=0." Inventory: equations (8.4.1)--(8.4.26); no notes, figures, tables,
% problems, references, or unresolved uncertainties.

\subsection{Electrodynamics in the Interaction Picture}

We now break up the Hamiltonian \eqref{eq:8.3.11} into a free-particle term $H_0$
and an interaction $H_{\mathrm{int}}$
\begin{align}
H&=H_0+H_{\mathrm{int}},
\tag{8.4.1}\label{eq:8.4.1}\\
H_0&=\int d^3x\left[
  \frac12\boldsymbol\Pi_{\!\perp}^{2}
  +\frac12(\boldsymbol\nabla\cdot\mathbf A)^2
  \right]+H_{\mathrm{matter},0},
\tag{8.4.2}\label{eq:8.4.2}\\
H_{\mathrm{int}}&=-\int d^3x\,\mathbf J\cdot\mathbf A
  +V_{\mathrm{Coul}}+H_{\mathrm{int,matter}},
\tag{8.4.3}\label{eq:8.4.3}
\end{align}
where $H_{\mathrm{matter},0}$ and $H_{\mathrm{int,matter}}$ are the
free-particle and interaction terms in $H_{\mathrm{matter}}$, and
$V_{\mathrm{Coul}}$ is the Coulomb interaction \eqref{eq:8.3.12}. The total Hamiltonian
\eqref{eq:8.4.1} is time-independent, so Eqs.~\eqref{eq:8.4.2} and \eqref{eq:8.4.3} can be evaluated at
any time we like (as long as both are evaluated at the same time), in
particular at $t=0$. As in Chapter 7, the transition to the interaction
picture is made by applying the similarity transformation
\begin{align}
H_I(t)
&=\exp(iH_0t)\,
  H_{\mathrm{int}}[\mathbf A,\boldsymbol\Pi_{\!\perp},Q,P]_{t=0}
  \exp(-iH_0t)
\notag\\
&=H_{\mathrm{int}}[\mathbf a(t),\boldsymbol\pi(t),q(t),p(t)],
\tag{8.4.4}\label{eq:8.4.4}
\end{align}
where $P$ here denotes the canonical conjugates to the matter fields $Q$, and
any operator $O(\mathbf x,t)$ in the interaction picture is related to its
value $O(\mathbf x,0)$ in the Heisenberg picture at $t=0$ by
\begin{equation}
O(\mathbf x,t)=\exp(iH_0t)O(\mathbf x,0)\exp(-iH_0t),
\tag{8.4.5}\label{eq:8.4.5}
\end{equation}
so that
\begin{equation}
i\dot O(\mathbf x,t)=[O(\mathbf x,t),H_0].
\tag{8.4.6}\label{eq:8.4.6}
\end{equation}
(We are dropping the subscript $\perp$ on $\boldsymbol\pi(\mathbf x)$.) Since
Eq.~\eqref{eq:8.4.5} is a similarity transformation, the equal-time commutation
relations are the same as in the Heisenberg picture:
\begin{align}
[a^i(\mathbf x,t),\pi^j(\mathbf y,t)]
&=i\left[\delta_{ij}\delta^3(\mathbf x-\mathbf y)
 +\frac{\partial^2}{\partial x^i\partial x^j}
  \frac{1}{4\pi\lvert\mathbf x-\mathbf y\rvert}\right],
\tag{8.4.7}\label{eq:8.4.7}\\
[a^i(\mathbf x,t),a^j(\mathbf y,t)]&=0,
\tag{8.4.8}\label{eq:8.4.8}\\
[\pi^i(\mathbf x,t),\pi^j(\mathbf y,t)]&=0,
\tag{8.4.9}\label{eq:8.4.9}
\end{align}
and likewise for the matter fields and their conjugates. For the same reason,
the constraints \eqref{eq:8.2.6} and \eqref{eq:8.3.8} still apply
\begin{align}
\boldsymbol\nabla\cdot\mathbf a&=0,
\tag{8.4.10}\label{eq:8.4.10}\\
\boldsymbol\nabla\cdot\boldsymbol\pi&=0.
\tag{8.4.11}\label{eq:8.4.11}
\end{align}
To establish the relation between $\boldsymbol\pi$ and $\dot{\mathbf a}$, we
must use Eq.~\eqref{eq:8.4.6} to evaluate $\dot{\mathbf a}$:
\begin{align*}
i\dot a_i(\mathbf x,t)
&=[a_i(\mathbf x,t),H_0]\\
&=i\int d^3y\left[
  \delta_{ij}\delta^3(\mathbf x-\mathbf y)
  +\frac{\partial^2}{\partial x^i\partial x^j}
   \frac{1}{4\pi\lvert\mathbf x-\mathbf y\rvert}
  \right]\pi_j(\mathbf y,t).
\end{align*}
We can replace $\partial/\partial x^j$ in the second term with
$-\partial/\partial y^j$, integrate by parts, and use Eq.~\eqref{eq:8.4.11}, yielding
\begin{equation}
\dot{\mathbf a}=\boldsymbol\pi
\tag{8.4.12}\label{eq:8.4.12}
\end{equation}
just as in the Heisenberg picture. The field equation is likewise determined
by
\begin{align*}
i\dot\pi_i(\mathbf x,t)
&=[\pi_i(\mathbf x,t),H_0]\\
&=-i\int d^3y\left[
  \delta_{ij}\delta^3(\mathbf x-\mathbf y)
  +\frac{\partial^2}{\partial x^i\partial x^j}
   \frac{1}{4\pi\lvert\mathbf x-\mathbf y\rvert}
  \right]
  \bigl(\boldsymbol\nabla\cdot\boldsymbol\nabla\cdot
        \mathbf a(\mathbf y,t)\bigr)_j,
\end{align*}
which (using Eqs.~\eqref{eq:8.4.10} and \eqref{eq:8.4.12}) just yields the usual wave equation
\begin{equation}
\Box\mathbf a=0.
\tag{8.4.13}\label{eq:8.4.13}
\end{equation}
Since $A^0$ is not an independent Heisenberg-picture field variable, but rather
a functional \eqref{eq:8.2.9} of the matter fields and their canonical conjugates that
vanishes in the limit of zero charges, we do not introduce any corresponding
operator $a^0$ in the interaction picture, but rather take
\begin{equation}
a^0=0.
\tag{8.4.14}\label{eq:8.4.14}
\end{equation}
The most general real solution of Eqs.~\eqref{eq:8.4.10}, \eqref{eq:8.4.13}, and \eqref{eq:8.4.14} may
be written
\begin{align}
a^\mu(x)=(2\pi)^{-3/2}\int\frac{d^3p}{\sqrt{2p^0}}\sum_\sigma
\Bigl[&e^{ip\cdot x}e^\mu(\mathbf p,\sigma)a_{\mathbf p,\sigma}
\notag\\[-2pt]
&+e^{-ip\cdot x}e^{\mu *}(\mathbf p,\sigma)
  a^\dagger_{\mathbf p,\sigma}\Bigr],
\tag{8.4.15}\label{eq:8.4.15}
\end{align}
where $p^0\equiv\lvert\mathbf p\rvert$; $e^\mu(\mathbf p,\sigma)$ are any two
independent `polarization vectors' satisfying
\begin{align}
\mathbf p\cdot\mathbf e(\mathbf p,\sigma)&=0,
\tag{8.4.16}\label{eq:8.4.16}\\
e^0(\mathbf p,\sigma)&=0,
\tag{8.4.17}\label{eq:8.4.17}
\end{align}
and $a_{\mathbf p,\sigma}$ are a pair of operator coefficients, with $\sigma$ a
two-valued index. By adjusting the normalization of $a_{\mathbf p,\sigma}$, we
can normalize the $e^\mu(\mathbf p,\sigma)$ so that the completeness relation
reads
\begin{equation}
\sum_\sigma e^i(\mathbf p,\sigma)e^j(\mathbf p,\sigma)^*
=\delta_{ij}-\frac{p_i p_j}{\lvert\mathbf p\rvert^2}.
\tag{8.4.18}\label{eq:8.4.18}
\end{equation}
For instance, we could take the $e(\mathbf p,\sigma)$ to be the same
polarization vectors that we encountered in Section 5.9:
\begin{equation}
e^\mu(\mathbf p,\mathord\pm1)=R(\widehat{\mathbf p})
\begin{bmatrix}
0\\[1pt]1/\sqrt2\\[1pt]\mathord\pm i/\sqrt2\\[1pt]0
\end{bmatrix},
\tag{8.4.19}\label{eq:8.4.19}
\end{equation}
where $R(\widehat{\mathbf p})$ is a standard rotation that carries the
three-axis into the direction of $\mathbf p$. Using Eqs.~\eqref{eq:8.4.18} and
\eqref{eq:8.4.12}, we can easily see that the commutation relations
\eqref{eq:8.4.7}--\eqref{eq:8.4.9} are satisfied if (and in fact only if) the operator
coefficients in Eq.~\eqref{eq:8.4.15} satisfy
\begin{align}
[a_{\mathbf p,\sigma},a^\dagger_{\mathbf p',\sigma'}]
&=\delta^3(\mathbf p-\mathbf p')\delta_{\sigma\sigma'},
\tag{8.4.20}\label{eq:8.4.20}\\
[a_{\mathbf p,\sigma},a_{\mathbf p',\sigma'}]&=0.
\tag{8.4.21}\label{eq:8.4.21}
\end{align}
As remarked before for massive particles, this result should be regarded not
so much as an alternative derivation of Eqs.~\eqref{eq:8.4.20} and \eqref{eq:8.4.21}, but rather
as a verification that Eq.~\eqref{eq:8.4.2} gives the correct Hamiltonian for free
massless particles of helicity $\mathord\pm1$. In the same spirit one can also
use Eqs.~\eqref{eq:8.4.12} and \eqref{eq:8.4.15} in Eq.~\eqref{eq:8.4.2} to calculate the free-photon
Hamiltonian
\begin{align}
H_0
&=\int d^3p\sum_\sigma\frac12p^0
  [a_{\mathbf p,\sigma},a^\dagger_{\mathbf p,\sigma}]_+
\notag\\
&=\int d^3p\sum_\sigma p^0
  \left(a^\dagger_{\mathbf p,\sigma}a_{\mathbf p,\sigma}
  +\frac12\delta^3(\mathbf p-\mathbf p)\right)
\tag{8.4.22}\label{eq:8.4.22}
\end{align}
which (aside from an inconsequential infinite $c$-number term) is just what we
should expect.

Finally, we record that the interaction \eqref{eq:8.4.4} in the interaction picture is
\begin{equation}
H_I(t)=-\int d^3x\,j_\mu(\mathbf x,t)a^\mu(\mathbf x,t)
 +V_{\mathrm{Coul}}(t)+H_{I,\mathrm{matter}}(t),
\tag{8.4.23}\label{eq:8.4.23}
\end{equation}
where in terms of the current $J$ in the Heisenberg picture
\begin{equation}
j_\mu(\mathbf x,t)\equiv
\exp(iH_0t)J_\mu(\mathbf x,0)\exp(-iH_0t),
\tag{8.4.24}\label{eq:8.4.24}
\end{equation}
while $V_{\mathrm{Coul}}(t)$ is the Coulomb term
\begin{align}
V_{\mathrm{Coul}}(t)
&=\exp(iH_0t)V_{\mathrm{Coul}}\exp(-iH_0t)
\notag\\
&=\frac12\int d^3x\,d^3y\,
  \frac{j^0(\mathbf x,t)j^0(\mathbf y,t)}
       {4\pi\lvert\mathbf x-\mathbf y\rvert}
\tag{8.4.25}\label{eq:8.4.25}
\end{align}
and $H_{I,\mathrm{matter}}(t)$ is the non-electromagnetic part of the matter
field interaction in the interaction picture:
\begin{equation}
H_{I,\mathrm{matter}}(t)=
\exp(iH_0t)H_{\mathrm{int,matter}}\exp(-iH_0t).
\tag{8.4.26}\label{eq:8.4.26}
\end{equation}
We have written $j_\mu a^\mu$ instead of $\mathbf j\cdot\mathbf a$ in Eq.
\eqref{eq:8.4.23}, but these are equal because $a^\mu$ has been defined to have
$a^0=0$.
